%% ma: L11-B16-0091
%% lop: 11
%% chuong: 5
%% bai: 16
%% dang: 11.16.4
%% loai: TL
%% muc: VD
%% dap_an: "-"
%% nguon: "Ảnh giáo viên gửi 25/11/2026: Giới hạn hàm số dạng 0/0 (3 bài, có ĐS), Bài 2 câu 9"
%% kien_thuc: "Lớp 11 Bài 16 (giới hạn của hàm số tại một điểm); phân tích đa thức thành nhân tử, hằng đẳng thức"
%% trang_thai: chinh-thuc
%% ghi_chu: "Đề gốc ghi mẫu x² - 3x - 2 và x² - 5x - 6: khi đó không có dạng vô định, giới hạn bằng -1/3, không khớp ĐS -2. Đã sửa thành x² - 3x + 2 và x² - 5x + 6 để khớp ĐS."
\begin{ex}
Tính giới hạn $A=\lim\limits_{x\to2}\left(\dfrac1{x^2-3x+2}+\dfrac1{x^2-5x+6}\right)$.
\loigiai{$\dfrac1{(x-1)(x-2)}+\dfrac1{(x-2)(x-3)}=\dfrac{(x-3)+(x-1)}{(x-1)(x-2)(x-3)}=\dfrac{2}{(x-1)(x-3)}$. Vậy $A=\dfrac{2}{1\cdot(-1)}=-2$.}
\end{ex}
